\chapter{Addition et soustraction}\label{ch:g3:addsub}

Ajouter, c'est mettre ensemble~; soustraire, c'est enlever, ou
trouver ce qui manque. Ce chapitre entraîne les deux techniques en
colonnes --- avec leurs célèbres retenues et emprunts --- et les
astuces mentales qui rendent souvent les colonnes inutiles.

\section{Addition en colonnes}

\begin{method}[Additionner en colonnes]\label{met:g3:addsub:add}
\begin{enumerate}
  \item Écrire les nombres l'un sous l'autre, unités sous unités,
        dizaines sous dizaines (aligner à \emph{droite})~;
  \item additionner les unités~; si la \omterm{def:g1:addition:def}{somme} dépasse $9$, écrire
        son chiffre des unités et \emph{retenir} la dizaine pour la
        colonne suivante (un petit $1$)~;
  \item continuer avec les dizaines, les centaines, \dots, en
        ajoutant la retenue à chaque fois.
\end{enumerate}
\end{method}

\begin{example}\label{ex:g3:addsub:add}
Calculer $457 + 368$~:
\[
\begin{array}{r}
{\scriptstyle 1}\ \ {\scriptstyle 1}\ \ \phantom{0} \\[-3pt]
4\,5\,7 \\
+\ 3\,6\,8 \\
\hline
8\,2\,5
\end{array}
\]
Unités~: $7 + 8 = 15$~: écrire $5$, retenir $1$. Dizaines~:
$5 + 6 + 1 = 12$~: écrire $2$, retenir $1$. Centaines~:
$4 + 3 + 1 = 8$. Résultat~: $825$.
\end{example}

\begin{example}[Addition mentale]\label{ex:g3:addsub:mental}
Chercher des nombres amis avant de se jeter sur les colonnes~:
\begin{itemize}
  \item \emph{faire une dizaine}~: $47 + 8 = 47 + 3 + 5 = 50 + 5 =
        55$~;
  \item \emph{sauter par parties}~: $256 + 130 = 256 + 100 + 30 =
        386$~;
  \item \emph{réordonner}~: $17 + 59 + 3 = (17 + 3) + 59 = 79$.
\end{itemize}
\end{example}

\begin{omfigure}
\begin{tikzpicture}[scale=1.05]
  \draw[-Stealth] (-0.2,0) -- (10.4,0);
  \foreach \x/\l in {0.56/{256}, 5.56/{356}, 7.06/{386}}
    { \draw (\x,-0.08) -- (\x,0.08);
      \node[below, font=\small] at (\x,-0.1) {$\l$}; }
  \draw[-Stealth, omDef, thick] (0.56,0.25) .. controls (2.8,1.15) ..
    (5.56,0.25);
  \node[omDef, font=\small] at (3.05,1.05) {$+100$};
  \draw[-Stealth, omThm, thick] (5.56,0.25) .. controls (6.3,0.8) ..
    (7.06,0.25);
  \node[omThm, font=\small] at (6.3,0.85) {$+30$};
\end{tikzpicture}

{\small L'addition mentale comme sauts sur la droite numérique~:
$256 + 130$ en deux bonds, d'abord les centaines, puis les dizaines.}
\end{omfigure}

\section{Soustraction en colonnes}

\begin{method}[Soustraire en colonnes]\label{met:g3:addsub:sub}
\begin{enumerate}
  \item Écrire le plus grand nombre en haut, aligné à droite~;
  \item soustraire colonne par colonne, en partant des unités~;
  \item si le chiffre du haut est trop petit, \emph{emprunter}~:
        prendre un à la place suivante sur la ligne du haut (il vaut
        dix dans la colonne en cours)~;
  \item \emph{vérifier}~: le résultat plus le nombre soustrait
        doit redonner le nombre du haut.
\end{enumerate}
\end{method}

\begin{example}\label{ex:g3:addsub:sub}
Calculer $603 - 147$~:
\[
\begin{array}{r}
6\,0\,3 \\
-\ 1\,4\,7 \\
\hline
4\,5\,6
\end{array}
\]
Unités~: $3 - 7$ impossible~; emprunter à travers~: $603$ devient
$5$ centaines, $9$ dizaines, $13$ unités. Puis $13 - 7 = 6$~;
dizaines $9 - 4 = 5$~; centaines $5 - 1 = 4$. Résultat~: $456$.
Vérification~: $456 + 147 = 603$.
\end{example}

\begin{omfigure}
\begin{tikzpicture}[scale=0.82]
  % before: 6 hundreds, 0 tens, 3 units
  \node[font=\small] at (2.5,2.25) {$603$ avant l'emprunt};
  \foreach \i in {0,...,5}
    \draw[omProp, fill=omProp!25] (0.62*\i,0.7) rectangle
      (0.62*\i+0.44,1.7);
  \node[below, font=\footnotesize] at (1.7,0.55) {$6$ centaines};
  \node[font=\footnotesize] at (4.35,1.2) {$0$ dizaine};
  \foreach \i in {0,1,2} \fill[omDef] (5.4+0.4*\i,1.2) circle (0.13);
  \node[below, font=\footnotesize] at (5.8,0.55) {$3$ unités};
  % after: 5 hundreds, 9 tens, 13 units
  \begin{scope}[yshift=-3.3cm]
  \node[font=\small] at (2.5,2.25) {$603$ après l'emprunt};
  \foreach \i in {0,...,4}
    \draw[omProp, fill=omProp!25] (0.62*\i,0.7) rectangle
      (0.62*\i+0.44,1.7);
  \node[below, font=\footnotesize] at (1.4,0.55) {$5$ centaines};
  \foreach \i in {0,...,8}
    \draw[omThm, very thick] (3.5+0.22*\i,0.7) -- (3.5+0.22*\i,1.7);
  \node[below, font=\footnotesize] at (4.4,0.55) {$9$ dizaines};
  \foreach \i in {0,...,6} \fill[omDef] (6.1+0.4*\i,1.45) circle (0.13);
  \foreach \i in {0,...,5} \fill[omDef] (6.1+0.4*\i,0.95) circle (0.13);
  \node[below, font=\footnotesize] at (7.3,0.55) {$13$ unités};
  \end{scope}
\end{tikzpicture}

{\small L'emprunt de $603 - 147$ en images~: une centaine est déliée
en dix dizaines, et une de ces dizaines en dix unités. Même
quantité, nouveaux habits --- $6\,\vert\,0\,\vert\,3$ devient
$5\,\vert\,9\,\vert\,13$, et chaque colonne peut maintenant
soustraire.}
\end{omfigure}

\begin{example}[La soustraction comme partie manquante]\label{ex:g3:addsub:missing}
\guillemotleft{} Lila a $35$ billes, Tom en a $52$~: combien de plus
a Tom~? \guillemotright{} est la soustraction $52 - 35$. De tête,
compter \emph{en montant} à partir de $35$~: $35 \to 40$ c'est
$5$, puis $40 \to 52$ c'est $12$~; en tout $5 + 12 = 17$. Compter en
montant est souvent plus facile qu'emprunter.
\end{example}

\begin{omfigure}
\begin{tikzpicture}[scale=1.05]
  \draw[-Stealth] (-0.2,0) -- (10.4,0);
  \foreach \x/\l in {1.0/{35}, 3.5/{40}, 9.5/{52}}
    { \draw (\x,-0.08) -- (\x,0.08);
      \node[below, font=\small] at (\x,-0.1) {$\l$}; }
  \draw[-Stealth, omDef, thick] (1.0,0.25) .. controls (2.2,0.9) ..
    (3.5,0.25);
  \node[omDef, font=\small] at (2.25,0.92) {$+5$};
  \draw[-Stealth, omThm, thick] (3.5,0.25) .. controls (6.5,1.2) ..
    (9.5,0.25);
  \node[omThm, font=\small] at (6.5,1.1) {$+12$};
\end{tikzpicture}

{\small Calculer $52 - 35$ en comptant en montant~: d'abord jusqu'au
$40$ ami, puis jusqu'à $52$. La réponse est le total des sauts~:
$5 + 12 = 17$.}
\end{omfigure}

\section{Ordres de grandeur}

\begin{method}[Estimer avant de calculer]\label{met:g3:addsub:estimate}
Avant tout calcul en colonnes, arrondir les nombres à des valeurs
amicales et calculer une \emph{estimation}. Après le calcul,
comparer~: un résultat loin de l'estimation signale une erreur.
\end{method}

\begin{example}
Pour $487 + 312$~: l'estimation est $500 + 300 = 800$. Le résultat
en colonnes, $799$, est proche~: plausible. Si l'on avait trouvé
$1\,099$ (un alignement raté), l'estimation aurait avertie.
\end{example}

\section{Exercices}

\begin{exercise}[$\star$]\label{exo:g3:addsub:1}
Calcule en colonnes~: $346 + 253$~; \quad $478 + 265$~; \quad
$1\,536 + 2\,876$.
\end{exercise}

\begin{exercise}[$\star$]\label{exo:g3:addsub:2}
Calcule en colonnes, puis vérifie par une addition~:
$587 - 234$~; \quad $805 - 348$~; \quad $1\,000 - 463$.
\end{exercise}

\begin{exercise}[$\star$]\label{exo:g3:addsub:3}
Calcule de tête en faisant des dizaines~: $38 + 7$~; \quad
$56 + 9$~; \quad $45 + 28$ (ajoute $30$, enlève $2$).
\end{exercise}

\begin{exercise}[$\star$]\label{exo:g3:addsub:4}
Calcule de tête par sauts~: $340 + 220$~; \quad $675 + 210$~; \quad
$512 - 300$~; \quad $840 - 220$.
\end{exercise}

\begin{exercise}[$\star$]\label{exo:g3:addsub:5}
Calcule $63 - 47$ en comptant en montant (comme dans
\cref{ex:g3:addsub:missing}), puis vérifie en colonnes.
\end{exercise}

\begin{exercise}[$\star$]\label{exo:g3:addsub:6}
Estime d'abord (arrondis aux centaines), puis calcule exactement~:
$492 + 218$~; \quad $913 - 388$.
\end{exercise}

\begin{exercise}[$\star$]\label{exo:g3:addsub:7}
Une école a $348$ filles et $296$ garçons. Combien d'élèves
en tout~? Écris l'opération, calcule, et réponds par une phrase.
\end{exercise}

\begin{exercise}[$\star$]\label{exo:g3:addsub:8}
Un livre a $224$ pages~; Nora en a lu $178$. Combien de pages
restent à lire~?
\end{exercise}

\begin{exercise}[$\star$]\label{exo:g3:addsub:9}
Complète les trous (travaille colonne par colonne)~:
\[
\begin{array}{r}
3\,?\,7 \\
+\ ?\,2\,? \\
\hline
8\,6\,2
\end{array}
\]
Y a-t-il une seule façon de remplir les trois trous~? Explique.
\end{exercise}

\begin{exercise}[$\star\star$]\label{exo:g3:addsub:10}
Une boulangerie a vendu $145$ croissants le matin et $89$
l'après-midi~; $17$ croissants restaient invendus à la fermeture.
Combien de croissants ont été cuits ce jour-là~? (Deux
étapes.)
\end{exercise}

\begin{exercise}[$\star\star$]\label{exo:g3:addsub:11}
Théo dit~: \guillemotleft{} pour calculer $500 - 199$, je calcule
$500 - 200$ puis j'ajoute $1$. \guillemotright{} A-t-il raison~?
Calcule des deux façons, et explique son astuce.
\end{exercise}
